\chapter{Approfondissements scientifiques}

\section{Role du chapitre}
Ce chapitre concentre l'analyse methodologique du memoire. Il ne repete pas les resultats du chapitre precedent et ne remplace pas la discussion generale. Son role est de preciser ce que les experiences permettent de conclure, ce qu'elles ne permettent pas de conclure, et comment les termes ambitieux du titre doivent etre compris.

\section{Lecture methodologique des scores supervises}
Les scores supervises doivent toujours etre lus avec leur protocole. Les protocoles exploratoires donnent des valeurs elevees sur Linux/auth, CICIDS2017 et CIC-DDoS2019. Les holdouts plus stricts donnent cependant F1 = 0,826251, 0,157827 et 0,997577. Le resultat CICIDS2017 est le plus instructif : la comparaison controlee montre qu'a features, hyperparametres, tailles et seeds apparies, le F1 passe de 0,999260 en split aleatoire stratifie a 0,157827 sur cinq holdouts fichier/scenario associes aux graines 42--46.

Cette experience ne signifie pas que CICIDS2017 est inutilisable ni qu'un autre modele echouerait necessairement. Elle montre que le RandomForest utilise dans cette configuration apprend fortement des signatures propres aux fichiers ou scenarios vus pendant l'entrainement. L'analyse par fichier precise le phenomene : DDoS reste partiellement detecte, PortScan presque pas, et Bot, Infiltration et WebAttacks ne sont pas retrouves dans les echantillons de test. L'ecart-type eleve ne doit donc pas etre interprete comme une simple variabilite due aux seeds ; il reflete surtout l'heterogeneite des scenarios tenus hors entrainement. Pour cette raison, le score exploratoire de 0,997163 n'est jamais utilise seul comme resultat principal.

Une passe complementaire compare plusieurs candidats legers sur les memes holdouts : RandomForest, ExtraTrees, HistGradientBoosting, LogisticRegression et SGDLogistic. Le meilleur F1 moyen est obtenu par LogisticRegression avec 0,233670, contre 0,157827 pour RandomForest. Ce gain est utile pour ameliorer le resultat strict, mais il ne suffit pas a conclure que le probleme est resolu : certains scenarios restent presque ou totalement non detectes. La lecture retenue est donc que le choix du modele compte, mais que la dependance au changement de distribution demeure le facteur dominant.

Le cas CIC-DDoS2019 est traite separement. Le score 0,999965 est conserve comme resultat exploratoire historique rattache a une famille DDoS, mais la preuve principale est le holdout strict sur cinq graines documente dans \texttt{random\_forest\_cic\_ddos2019\_metrics.csv}. Le dataset UNSW-NB15 n'est pas fusionne avec CIC-DDoS2019 ; il apparait seulement dans l'ablation controlee distincte.

\section{HDFS/BGL et risque de fuite}
Les premiers essais HDFS/BGL ont une valeur d'integration, mais ils ne constituent pas une validation robuste lorsque le seuil ou la contamination depend du jeu evalue. Utiliser implicitement l'information du test pour regler un hyperparametre peut surestimer fortement le F1, en particulier sur BGL. Le memoire ne presente donc pas ces scores exploratoires comme preuve principale.

La correction apportee est la branche Drain3 train--test : extraction de templates, fenetres temporelles, separation des ensembles et evaluation sur test. Cette version rend le protocole plus coherent avec la litterature sur les logs sequentiels. Elle reste toutefois limitee : HDFS demeure difficile, BGL reste tres favorable sur le split local, et une validation plus forte demanderait repetitions de splits, seuil fixe sur validation et comparaison a des modeles sequentiels specialises.

\section{Agents, parsing et portee du titre}
Le terme agent intelligent est employe dans un sens operationnel. Les agents disposent de capacites declarees, d'une memoire locale, d'un heartbeat, d'une politique de choix et de handlers multiples. La memoire adaptative ajoute une capitalisation persistante de certains historiques et retours analyste, mais elle ne transforme pas les agents en agents cognitifs autonomes capables de planification generale ou d'apprentissage par renforcement. La contribution se situe donc dans l'orchestration logicielle auditable et l'adaptation comportementale controlee, pas dans une nouvelle architecture agentique d'IA generale.

Le parsing doit etre compris de la meme maniere. Le prototype sait reconnaitre plusieurs familles controlees, extraire des champs utiles et conserver les lignes inconnues ou incompletes. Il ne demontre pas une adaptation autonome illimitee a tout format nouveau. Sa robustesse vient surtout de la degradation controlee : ne pas perdre silencieusement l'information lorsque le format sort du cadre connu.

La distribution est egalement locale. La campagne Redis de six heures montre une repartition multi-processus et une reprise de taches non acquittees. Une comparaison controlee monolithique-versus-agents a ete ajoutee sur les memes taches et le meme volume : elle ne montre pas de gain de debit sur petit volume, mais elle quantifie l'overhead d'orchestration et confirme la reprise d'une tache simulee comme non acquittee. La modularite et la tracabilite sont soutenues par la conception logicielle et les mecanismes d'audit ; la resilience locale dispose, elle, d'une preuve experimentale directe par panne simulee.

\section{Synthese des risques de validite}
\begin{table}[H]
\centering
\caption{Risques methodologiques et formulation retenue}
\label{tab:risques-methodologiques}
\scriptsize
\begin{tabularx}{\textwidth}{p{3.2cm}p{4.2cm}X}
\toprule
\textbf{Point critique} & \textbf{Risque} & \textbf{Formulation retenue} \\
\midrule
Scores supervises & Confusion entre score exploratoire et generalisation. & Les scores sont associes a leur protocole ; CICIDS2017 est accompagne du holdout controle. \\
CIC-DDoS2019/UNSW & Attribution dataset ambigue. & CIC-DDoS2019 et UNSW-NB15 sont separes ; le resultat principal DDoS est le holdout strict. \\
HDFS/BGL & Fuite par seuil ou contamination. & Les premiers scores sont exploratoires ; Drain3 train--test devient la lecture principale. \\
Dashboard & Confusion entre validation fonctionnelle et evaluation utilisateur. & Les captures prouvent le fonctionnement ; une user study reste hors perimetre. \\
Wazuh/Fail2ban & Comparaison abusive avec des outils matures. & \logminer{} est presente comme couche analytique complementaire, pas comme substitut. \\
Agents intelligents & Sur-promesse sur l'autonomie ou le debit. & Agents logiciels multi-taches avec orchestration ; comparaison controlee mesurant surtout overhead et reprise, pas acceleration brute. \\
Memoire adaptative & Confusion entre feedback conserve, reduction mesuree des faux positifs et reentrainement automatique. & Memoire persistante auditable pour priorisation et contexte ; reduction reelle des faux positifs et apprentissage adaptatif complet non encore demontres. \\
Scalabilite & Benchmarks locaux trop courts pour une preuve industrielle. & Les mesures documentent la faisabilite locale ; la scalabilite multi-machine reste une perspective. \\
\bottomrule
\end{tabularx}
\end{table}

\section{Exploitation scientifique prioritaire}
Les resultats les plus utiles pour la suite sont ceux qui isolent une question scientifique claire. La comparaison controlee CICIDS2017 montre l'effet du protocole de separation. L'ablation global-versus-familles sur cinq graines montre que le routage ne domine pas automatiquement un modele global en espace commun. La branche Drain3 montre que les logs sequentiels demandent une representation plus riche que la ligne brute.

Les prolongements doivent donc suivre trois axes : etendre les comparaisons controlees a d'autres familles, tester des splits temporels ou multi-environnements, et transformer la memoire adaptative actuelle en boucle de reentrainement controlee fondee sur le feedback analyste du dashboard. Cette boucle devra surveiller le concept drift, comparer tout nouveau modele a un modele courant sur un jeu de validation gele, puis conserver l'ancien modele pour retour arriere.
